\section{Even degrees prime to three}\label{sec:even}

In this section and the next, $M$ is even and prime to $3$, and $A$ is a
sextuple at level $M$ that contains no pair, is not split and generates
$\ZZ/M$. We look for a chain of moves that succeeds, as required by
\Cref{prop:induction}. The definitions of \Cref{sec:odd} (maximal orders, the
embedding $\kappa$, the function $F$, moves $(\ell,x)$ for $A$) and
\Cref{lem:movecrit,lem:charrel,lem:fiberline} do not use that $M$ is odd, and
we use them as they stand.

The \emph{odd part} of a positive integer $f$ is $f/2^{v_{2}(f)}$. Let
$e^{*}$ be an odd part of an order of an element of $A$ that is maximal for
divisibility among these odd parts, and let $2^{J}e^{*}$ be the largest order
of an element of $A$ with odd part $e^{*}$. Then $2^{J}e^{*}$ is a maximal
order: an order that is a proper multiple of it would have an odd part
divisible by $e^{*}$, hence equal to $e^{*}$, and a larger power of $2$. This
section treats $J=1$, and \Cref{sec:four} treats $J\ge2$.

\subsection{Every maximal order is even}

\begin{lemma}\label{lem:maxeven}
If $A$ has no move, then every maximal order of $A$ is even. In particular
$J\ge1$.
\end{lemma}

\begin{proof}
Let $e$ be an odd maximal order. Since $3\nmid M$, $e$ is prime to $6$, and
$e>1$. As in \Cref{sec:odd}, $F$ is odd, $\|F\|=2N_{e}\le12$, $F(z_{0})\ge1$
when $\kappa(z_{0})\in A$ has order $e$, and $F$ is split along
$(K_{p})_{p\mid e}$ on every coset of $H$, so \Cref{prop:core} applies. In its
case (1), $F\ge1$ on a line $wK_{p}$ with $p\in\{5,7\}$; as $F$ is odd,
$-1\notin K_{p}$, so $e\ne p$, and $|K_{p}|\ge4\ge(p+1)/2$, so $A$ has a move
by \Cref{lem:fiberline}. In case (2), $\|F\|=12$, so the six elements of $A$
have order $35$ and lie in the subgroup of order $35$; they do not generate
$\ZZ/M$, as $M$ is even.
\end{proof}

\subsection{The top order}\label{ssec:top}
From now on in this section $J=1$, and we put $e=2e^{*}$. If $e^{*}=1$, every
odd part is $1$, the largest order is $2$, and $A$ consists of six copies of
$M/2$, which contains a pair; so $e^{*}>1$, and $e^{*}$ is prime to $6$. Let
$U^{*}=(\ZZ/e^{*})^{\times}$. For $v\in U^{*}$ let
\[
  x_{v}=\tfrac{M}{e}\,u\quad(u\ \text{odd},\ u\equiv v\bmod e^{*}),\qquad
  y_{v}=x_{v}+\tfrac M2=\tfrac{M}{e^{*}}\cdot\tfrac{v}{2},
\]
where $v/2$ is taken modulo $e^{*}$. The $x_{v}$ are the elements of order $e$
and the $y_{v}$ those of order $e^{*}$, and
\[
  -x_{v}=x_{-v},\quad -y_{v}=y_{-v},\quad 2x_{v}=y_{2v},\quad
  -2x_{v}=y_{-2v},\quad T(x_{v})=\{x_{v},\,y_{v},\,y_{-2v},\,M/2\}.
\]
Reduction modulo $e^{*}$ identifies $U=(\ZZ/e)^{\times}$ with $U^{*}$, and
$\kappa(u)=x_{u\bmod e^{*}}$. So the function $F$ of \Cref{sec:odd} becomes
$F(v)=\mult_{A}(x_{v})-\mult_{A}(x_{-v})$, and for $p\mid e^{*}$ the subgroup
$K_{p}$ becomes the group of units congruent to $1$ modulo $e^{*}/p$; we
write $H^{*}$ for the product of these subgroups, which is direct
(\Cref{lem:Kp}). Put
\[
  E(v)=\mult_{A}(y_{v})-\mult_{A}(y_{-v}).
\]
Both $F$ and $E$ are odd, and since $A$ contains no pair,
$|F(v)|=\mult_{A}(x_{v})+\mult_{A}(x_{-v})$ and likewise for $E$; so
$\|F\|=2N_{e}$ and $\|E\|=2N_{e^{*}}$, where $N_{f}$ is the number of
elements of $A$ of order $f$. The elements of $A$ of order $e$ are the
\emph{tops}; we say that $v$ is a top if $x_{v}\in A$, that is, if
$F(v)\ge1$. The elements of order other than $e$ and $e^{*}$ are \emph{low},
and there are $6-N_{e}-N_{e^{*}}$ of them.

\begin{lemma}\label{lem:topfibres}
Let $q\mid e^{*}$ be prime and $v\in U^{*}$. The elements $x_{w}$, $w\in
vK_{q}$, lie in $\Phi_{q}(x_{v})$, and so do the elements $y_{w}$, $w\in
vK_{q}$, if $e^{*}\ne q$. If $A$ contains $(q+1)/2$ distinct elements of one
of these two kinds, then $A$ has a move.
\end{lemma}

\begin{proof}
For $w\in vK_{q}$ we have $w\equiv v$ modulo $e^{*}/q$. The representatives
$u,u'$ of $v,w$ in the definition of $x$ can be chosen odd and congruent
modulo $2e^{*}/q$, so $x_{w}-x_{v}\in(M/e)(e/q)\ZZ=(M/q)\ZZ$; likewise
$y_{w}-y_{v}=(M/e^{*})(w-v)/2\in(M/q)\ZZ$. The last assertion is
\Cref{lem:movecrit}, applied with $\ell=q$, which differs from the orders $e$
and $e^{*}$ of these elements.
\end{proof}

\begin{lemma}\label{lem:E1}
Let $E_{1}(w)=E(w)-F(w)+F(w/2)$. Then $E_{1}$ is odd and split along
$(K_{p})_{p\mid e^{*}}$ on every coset of $H^{*}$ in $U^{*}$. It is the
function $w\mapsto\mult_{A_{1}}(y_{w})-\mult_{A_{1}}(y_{-w})$ of the multiset
\[
  A_{1}=A-\sum_{v\ \text{top}}\{x_{v}\}+\sum_{v\ \text{top}}\{y_{-v},\,y_{2v},
  \,M/2\},
\]
the sums running over the tops with their multiplicities, and $A_{1}$ is
the result of the chain of moves $(T(x_{v}),\{x_{v}\})$ over the tops.
\end{lemma}

\begin{proof}
Let $\chi$ be an odd primitive character modulo $e^{*}$. The elements $a$ of
$A$ whose order is divisible by $e^{*}$ have odd part $e^{*}$, by maximality,
so they are the $x_{v}$ and $y_{v}$ in $A$. For $y_{v}$ we have
$z_{a}=v/2$, and for $x_{v}$ we have $z_{a}\equiv v$ modulo $e^{*}$ and the
factor $1-\chi(2)$; and $\varphi(e)=\varphi(e^{*})$. Since $\chi$ is odd,
$\sum_{v}\mult_{A}(y_{v})\bar\chi(v)=\tfrac12\hat E(\chi)$ with $\hat
f(\chi)=\sum_{v}f(v)\bar\chi(v)$, and likewise for $F$. So \eqref{eq:charrel}
reads $\chi(2)\hat E(\chi)+(1-\chi(2))\hat F(\chi)=0$, that is, $\hat
E(\chi)=(1-\bar\chi(2))\hat F(\chi)$. As $\widehat{F(\cdot/2)}(\chi)=
\bar\chi(2)\hat F(\chi)$, we get $\hat E_{1}(\chi)=0$. For even $\chi$,
$\hat E_{1}(\chi)=0$ because $E_{1}$ is odd. So $\hat E_{1}$ vanishes at
every primitive character, and the proof of \Cref{lem:splitF}, with
\Cref{lem:Kp}(c) for $e^{*}$, shows that $E_{1}$ is split on every coset of
$H^{*}$.

The move $(T(x_{v}),\{x_{v}\})$ removes $x_{v}$ and adds
$-y_{v}=y_{-v}$, $-y_{-2v}=y_{2v}$ and $-M/2=M/2$. The function of
$A_{1}$ at $w$ receives $E(w)$ from $A$, $+1$ for each top $v=-w$ and $-1$ for
each top $v=w$, which is $-F(w)$, and $+1$ for each top $v=w/2$ and $-1$ for
each top $v=-w/2$, which is $F(w/2)$.
\end{proof}

\subsection{Elementary moves}

\begin{lemma}\label{lem:elementary}
Let $v$ be a top and $x=x_{v}$. Each of the following is a move that
lowers the profile, or leaves at most four elements after the removal of
pairs.
\begin{enumerate}[label=\textup{(\alph*)}]
\item $y_{v}=x+M/2\in A$: with $S=\{x,x+M/2\}\subseteq T(x)$,
$A''=A-x-(x+M/2)+\{2x,\,M/2\}$.
\item $y_{-2v}=-2x\in A$: with $S=\{x,-2x\}\subseteq T(x)$,
$A''=A-x-(-2x)+\{y_{-v},\,M/2\}$.
\item $2v$ is also a top: with $\sigma=T(x_{2v})+T(x_{v})$ and
$S=\{x_{2v},x_{v}\}$, $A''=A-x_{2v}-x_{v}+\{y_{4v},\,y_{-v}\}$.
\item $M/2\in A$: with $S=\{x,M/2\}\subseteq T(x)$,
$A''=A-x-M/2+\{y_{-v},\,y_{2v}\}$.
\end{enumerate}
\end{lemma}

\begin{proof}
These are instances of \Cref{def:move}. In (c), $\sigma$ is algebraic by
\Cref{lem:algtools}(b),(e), $\sigma-S=\{y_{2v},y_{-4v},M/2,y_{v},y_{-2v},
M/2\}$, and $-(\sigma-S)=\{y_{-2v},y_{4v},M/2,y_{-v},y_{2v},M/2\}$; after
removing the pairs $\{y_{2v},y_{-2v}\}$ and $\{M/2,M/2\}$ the stated multiset
remains. In each case one element of order
$e$ is removed and only elements of order $e^{*}$ or $2$ are added, while the
elements of order larger than $e$ stay. So if six elements remain they form a
sextuple of smaller profile.
\end{proof}

From now on we assume that $A$ has none of the moves of
\Cref{lem:elementary,lem:topfibres}.

\begin{lemma}\label{lem:signs}
For every $w\in U^{*}$ the numbers $E(w)$, $-F(w)$ and $F(w/2)$ do not have
opposite signs, so
\[
  |E_{1}(w)|=|E(w)|+|F(w)|+|F(w/2)|.
\]
Consequently:
\begin{enumerate}[label=\textup{(\roman*)}]
\item $\|E_{1}\|=2N_{e^{*}}+4N_{e}\le24$;
\item at every top $v$, $E_{1}(v)\le-1$ and $E_{1}(2v)\ge1$, so $E_{1}\ne0$;
\item the elements of order $e^{*}$ of $A_{1}$ contain no pair, and
$\mult_{A_{1}}(y_{w})$ is the positive part of $E_{1}(w)$;
\item the negative part of $E_{1}$ is the counting function of the multiset
$\{w:\ y_{-w}\in A\}+\{\text{tops}\}+\{-2v:\ v\ \text{top}\}$.
\end{enumerate}
\end{lemma}

\begin{proof}
If $E(w)\ge1$ and $F(w)\ge1$, then $y_{w}=x_{w}+M/2$ and $x_{w}$ lie in
$A$, which is \Cref{lem:elementary}(a); the same with $-w$ if both are
negative. If $F(w)\le-1$ and $F(w/2)\le-1$, then $-w$ and $-w/2$ are tops,
which is (c); the same with $w$ if $F(w)\ge1$ and $F(w/2)\ge1$. If
$E(w)\ge1$ and $F(w/2)\le-1$, then $y_{w}$ and $x_{-w/2}$ lie in $A$, and
$y_{w}=-2x_{-w/2}$, which is (b); the same with $-w$ in the other case. So
the three numbers never have opposite signs, and the displayed identity
follows. Summing it gives (i), with $N_{e}+N_{e^{*}}\le6$. At a top $v$,
$-F(v)\le-1$ and $F(2v/2)\ge1$, which gives (ii), and $N_{e}\ge1$ since $e$
is an order of $A$. The multiplicity of $y_{w}$ in $A_{1}$ is the sum of the
positive parts of the three numbers, and that of $y_{-w}$ the sum of their
negative parts; so one of them vanishes, which is (iii), and (iv) is the
same computation for the negative parts.
\end{proof}

\subsection{The compound move}
A \emph{non-cancelling line sum} is a function $f=\sum_{i}s_{i}\mathbf
1_{L_{i}}$, for a finite family of lines $L_{i}$ (repetitions allowed) and
signs $s_{i}\in\{\pm1\}$, such that $|f|=\sum_{i}\mathbf 1_{L_{i}}$.

\begin{lemma}\label{lem:linesum}
If $E_{1}$ is a non-cancelling line sum of lines along the subgroups $K_{p}$,
$p\mid e^{*}$, and $e^{*}$ is not prime, then $A$ has a chain of moves that
is direct or lowering.
\end{lemma}

\begin{proof}
Pass from $A$ to $A_{1}$ (\Cref{lem:E1}). By \Cref{lem:signs}(iii), the
elements of order $e^{*}$ of $A_{1}$ are the $y_{w}$ for $w$ in the lines
$L_{i}$ with $s_{i}=1$, each counted once per line. Let $L=wK_{p}$ be such a
line and $y=y_{w}$. By \Cref{lem:topfibres} the $y_{w'}$, $w'\in L$, are
$|K_{p}|$ distinct elements of order $e^{*}$ of the fibre $\Phi_{p}(y)$. If
$p^{2}\mid e^{*}$, they are the whole fibre; if $p\parallel e^{*}$, the fibre
consists of them and one element $y_{0}$ of order $e^{*}/p$. Since $e^{*}$ is
not prime, $py\ne0$, and the move $(\sigma_{p,y},S)$ with $S=\{y_{w'}:w'\in
L\}$ removes these elements and adds $py$, and $-y_{0}$ if $p\parallel e^{*}$:
at most two elements of order dividing $e^{*}/p$. We do this for each positive
line in turn.

Each line has $|K_{p}|\ge4$ points, and the positive lines carry
$\|E_{1}\|/2=N_{e^{*}}+2N_{e}$ points in all, so there are at most
$(N_{e^{*}}+2N_{e})/4$ of them. The multiset $A_{1}$ contains $N_{e}$ copies
of $M/2$, which is not in $A$ by \Cref{lem:elementary}(d); after removing pairs,
$\epsilon\in\{0,1\}$ copies remain, with $\epsilon\equiv N_{e}$ modulo $2$.
The final multiset therefore has at most
\[
  (6-N_{e}-N_{e^{*}})+\epsilon+2\cdot\frac{N_{e^{*}}+2N_{e}}{4}
  =6+\epsilon-\frac{N_{e^{*}}}{2}\le7
\]
elements. It is balanced (\Cref{lem:move}), so its size is even and at most
$6$. It contains the low elements of $A$, at most one copy of $M/2$ and
elements of order smaller than $e^{*}$; so it has no element of order $e$,
the same elements of order larger than $e$ as $A$, and the chain is direct or
lowering.
\end{proof}

\subsection{Split functions of mass at most eight}

\begin{lemma}\label{lem:S}
Let $K_{1},\dots,K_{r}$ be subgroups of a finite abelian group whose product
$H$ is direct, with pairwise distinct orders $|K_{i}|\ge4$, such that the
orders $4$ and $5$ do not both occur. Let $C$ be a coset of $H$ and $f\colon
C\to\ZZ$ split along $(K_{1},\dots,K_{r})$, not identically zero, with
$\sum_{C}|f|\le8$. Then $f$ is one of
\begin{enumerate}[label=\textup{(\roman*)}]
\item $\sigma\mathbf 1_{L}$, with $\sigma=\pm1$ and $L$ a line along some
$K_{i}$;
\item $2\sigma\mathbf 1_{L}$, with $L$ a line along the $K_{i}$ of order $4$;
\item $\sigma\mathbf 1_{L}+\sigma'\mathbf 1_{L'}$, with $L\ne L'$ lines along
the $K_{i}$ of order $4$;
\item a \emph{cross} $\sigma(\mathbf 1_{L}-\mathbf 1_{L'})$, with $L$ a line
along the $K_{i}$ of order $6$, $L'$ a line along the $K_{j}$ of order $4$, and
$|L\cap L'|=1$.
\end{enumerate}
In particular $\sum_{C}|f|\ge\min_{i}|K_{i}|$.
\end{lemma}

\begin{proof}
By induction on $r$; for $r=0$ there is nothing to prove. For $r=1$, $f$ is
constant on the line $C$ by (S4), with a value $c\ne0$ and $|c|\,|K_{1}|\le8$;
so $|c|=1$, case (i), or $|c|=2$ and $|K_{1}|=4$, case (ii). Let $r\ge2$, and
number the subgroups so that $K=K_{1}$ has the largest order; then $|K|\ge5$
and $|K|>m'=\min_{i\ge2}|K_{i}|\ge4$. We use the slices $S_{c}$, $c\in K$, of
\Cref{sec:primetosix}.

\emph{Case A: $f$ is $K$-invariant.} Then $f$ is constant on the lines along
$K$, a nonzero value $c$ gives $|c|\,|K|\le8$, so $|c|=1$, and a second
line on which $f\ne0$ would add $|K|\ge5$ to the mass. This is (i).

\emph{Case B1: $\mu(c_{0})=0$ for some $c_{0}$.} By (S2), $f$ is split along
$(K_{2},\dots,K_{r})$ on every slice. By induction a nonzero slice carries
mass at least $m'\ge4$, so at most two slices are nonzero. If one is, $f$ is
one of (i)--(iv) inside it. If two are, each carries mass $4$ and is of the
form $\sigma\mathbf 1_{L}$ with $|L|=4$; this is (iii).

\emph{Case B2: $f$ is not $K$-invariant and $\mu\ge1$ on every slice.} As
$\sum_{c}\mu(c)\le8<2|K|$, some $c_{s}$ has $\mu(c_{s})=1$. For $c\in K$ the
function $g_{c}(z)=f(z)-f(zc^{-1}c_{s})$ on $S_{c}$ is split along
$(K_{2},\dots,K_{r})$ by (S1), with mass at most $\mu(c)+1\le8$. If every
$g_{c}$ vanished, $f$ would be $K$-invariant. Let $D$ be the set of $c$ with
$g_{c}\ne0$, and $d=|D|\ge1$. For $c\in D$ the induction gives $\mu(c)+1\ge
m'$. Hence
\[
  8\ge\sum_{c}\mu(c)\ge d(m'-1)+(|K|-d),\qquad\text{so}\qquad |K|+d(m'-2)\le8.
\]
With $|K|\ge m'+1$ and $d\ge1$ this forces $m'=4$, $d=1$ and $|K|\le6$; the
value $|K|=5$ is excluded, as the orders $4$ and $5$ do not both occur. So
$|K|=6$, all the inequalities are equalities, $\mu(c_{1})=3$ for the element
$c_{1}$ of $D$, and $\mu=1$ on the other five slices. Off $S_{c_{1}}$, $f$ is
the translate of its restriction to $S_{c_{s}}$, so it equals $\sigma$ at one
point of each slice, and these points lie on one line $L$ along $K$. On
$S_{c_{1}}$, $g_{c_{1}}$ has mass $4$, so by induction it is $\tau\mathbf
1_{L'}$ for a line $L'$ along the subgroup of order $4$; then $f=\sigma\delta_{z_{1}}
+\tau\mathbf 1_{L'}$ on $S_{c_{1}}$, where $z_{1}$ is the point of $L$ in
$S_{c_{1}}$, and the mass $3$ forces $z_{1}\in L'$ and $\tau=-\sigma$. This
is (iv). The last assertion holds in each case.
\end{proof}

For $e^{*}$ prime to $6$ the subgroups $K_{p}$, $p\mid e^{*}$, satisfy the
hypotheses of \Cref{lem:S}: by \Cref{lem:Kp}(a), $|K_{p}|$ is $p-1$ or $p$,
these pairs of integers are disjoint for different primes $p\ge5$, and the
order $4$ occurs only for $5\parallel e^{*}$, the order $5$ only for
$25\mid e^{*}$. A cross needs $5\parallel e^{*}$ and $7\parallel e^{*}$. The
functions (i)--(iii) are non-cancelling line sums, and a cross is not.

\subsection{The cases}
We now show that $E_{1}$ is a non-cancelling line sum, except when $e^{*}$ is
$35$ or $55$, which we treat afterwards.

\begin{lemma}\label{lem:eprime}
$e^{*}$ is not prime.
\end{lemma}

\begin{proof}
If $e^{*}=p$, then $K_{p}=U^{*}=H^{*}$, and $E_{1}$ is constant by (S4) and
odd, so $E_{1}=0$, against \Cref{lem:signs}(ii).
\end{proof}

\begin{proposition}\label{prop:nonsquarefree}
If $e^{*}$ is not squarefree, then $E_{1}$ is a non-cancelling line sum.
\end{proposition}

\begin{proof}
Let $R=\prod_{p^{2}\mid e^{*}}p^{v_{p}(e^{*})-1}$; then $R\ge5$, since
$e^{*}$ is prime to $6$. Every $K_{p}$ consists of units congruent to $1$
modulo $R$, as $R\mid e^{*}/p$, and the group of such units has order
$\prod_{p\parallel e^{*}}(p-1)\prod_{p^{2}\mid e^{*}}p=|H^{*}|$; so $H^{*}$ is
the group of units congruent to $1$ modulo $R$. Hence $-1,2,-2,4\notin H^{*}$,
since $R$ divides none of $2,1,3,3$, and $-4\in H^{*}$ only if $R=5$.

Let $v$ be a top. The cosets $X_{1}=vH^{*}$, $X_{2}=-vH^{*}$, $X_{3}=2vH^{*}$,
$X_{4}=-2vH^{*}$ are distinct, and $E_{1}\ne0$ on each of them by
\Cref{lem:signs}(ii), so each carries mass at least $4$ (\Cref{lem:S}). By
oddness $\mu(X_{1})=\mu(X_{2})=a$ and $\mu(X_{3})=\mu(X_{4})=b$, and
$2a+2b\le24$; so $a,b\le8$, and every other coset, which has the same mass as
its negative, carries at most $4$. So on every coset $E_{1}$ is zero or one of
the functions of \Cref{lem:S}. If no coset carries a cross, $E_{1}$ is a
non-cancelling line sum.

Suppose a coset carries a cross. Then $5\parallel e^{*}$ and $7\parallel
e^{*}$, so the primes of $R$ are at least $11$, $R\ne5$ and $-4\notin H^{*}$.
The coset has mass $8$, so it is one of the $X_{i}$, and $\{a,b\}=\{4,8\}$.
Then $\|E_{1}\|=24=2N_{e^{*}}+4N_{e}$, so $N_{e}=6$, $N_{e^{*}}=0$, and $E_{1}$
vanishes off the $X_{i}$. By \Cref{lem:signs}(iv) the negative part of
$E_{1}$ counts the multiset $T+(-2T)$, $T$ the six tops. A top $t\in
X_{3}\cup X_{4}$ would have $E_{1}(2t)\ne0$ with $2t\in\pm4vH^{*}$, and this
coset is none of the $X_{i}$, as $\pm4,\pm2\notin H^{*}$, $\mp4\notin H^{*}$.
So $T\subseteq X_{1}\cup X_{2}$ and $-2T\subseteq X_{3}\cup X_{4}$, and the
negative mass of $E_{1}$ on $X_{1}\cup X_{2}$ is $|T|=6$. By oddness it is
the total mass $a$ of $E_{1}$ on $X_{1}$, which is $4$ or $8$; a
contradiction.
\end{proof}

\begin{proposition}\label{prop:threeprimes}
If $e^{*}$ is squarefree with at least three prime factors, then $E_{1}$ is
a non-cancelling line sum.
\end{proposition}

\begin{proof}
Here $H^{*}=U^{*}$. Let $p_{1}$ be the largest prime factor, so $p_{1}\ge11$,
and $K=K_{p_{1}}$, of order $p_{1}-1\ge10$. In the coordinates
$U^{*}=(\ZZ/p_{1})^{\times}\times(\ZZ/(e^{*}/p_{1}))^{\times}$ the group $K$ is
the first factor, the product $H'$ of the other $K_{p}$ is the second, and
the slices are $S_{c}=\{c\}\times(\ZZ/(e^{*}/p_{1}))^{\times}$. Write
$-1=\varepsilon\varepsilon'$ with $\varepsilon=(-1,1)\in K$ and
$\varepsilon'=(1,-1)\in H'$. Then $\mu(-c)=\mu(c)$, and for $z\in S_{c}$,
$E_{1}(z\varepsilon)=E_{1}(-z\varepsilon')=-E_{1}(z\varepsilon')$; so by
(S1) the function
\[
  g(z)=E_{1}(z)+E_{1}(z\varepsilon')=E_{1}(z)-E_{1}(z\varepsilon)
  \tag{$*$}
\]
is split along $(K_{p})_{p\ne p_{1}}$ on $S_{c}$, with mass at most
$2\mu(c)$, and $g(z\varepsilon')=g(z)$. The element $\varepsilon'$ lies in no
$K_{p}$ with $p\ne p_{1}$, since it is congruent to $-1$ modulo a third prime
factor. A top in $S_{c}$ makes the slices $S_{c},S_{-c},S_{2c},S_{-2c}$
nonzero (\Cref{lem:signs}(ii)), and they are distinct, as $p_{1}\ge11$.

\emph{Case Z: some slice vanishes.} By (S2), $E_{1}$ is split along
$(K_{p})_{p\ne p_{1}}$ on every slice, so nonzero slices carry mass at least
$4$ (\Cref{lem:S}). For a top in $S_{c}$, the four slices above carry
$2\mu(c)+2\mu(2c)\le24$, so $\mu(c),\mu(2c)\le8$ and every other slice
carries at most $4$. So \Cref{lem:S} applies on every slice, and without a
cross $E_{1}$ is a non-cancelling line sum. A cross has mass $8$, so for each
top it lies on one of the four slices of that top; then $\|E_{1}\|=24$,
$N_{e}=6$, $N_{e^{*}}=0$, and $E_{1}$ vanishes off these four slices, with
$\{\mu(c),\mu(2c)\}=\{4,8\}$. A top $t$ in $S_{\pm2c}$ would have $2t$ in
$S_{\pm4c}$, which is none of the four slices, as $4\not\equiv\pm1,\pm2$
modulo $p_{1}$; but $E_{1}(2t)\ne0$. So the six tops lie in $S_{c}\cup
S_{-c}$, where the negative mass of $E_{1}$ is then $6$ by
\Cref{lem:signs}(iv), while by oddness it is $\mu(c)\in\{4,8\}$.

\emph{Case NZ1: some $\mu(c)=1$.} In $(*)$, $g$ has mass at most $2<4$, so
$g=0$; then $E_{1}(z_{0}\varepsilon')=-E_{1}(z_{0})\ne0$ at the support point
$z_{0}$ of $S_{c}$, and $z_{0}\varepsilon'\ne z_{0}$ is a second point of
$S_{c}$ where $E_{1}\ne0$. A contradiction.

\emph{Case NZ2: all $\mu(c)\ge2$.} Then $2(p_{1}-1)\le24$, so $p_{1}$ is $11$
or $13$. On a slice of mass $2$, $g$ has mass at most $4$; if it is not
zero, it is $\pm\mathbf 1_{L}$ for a line $L$ along $K_{5}$ (\Cref{lem:S}),
and $g(z\varepsilon')=g(z)$ gives $L\varepsilon'=L$, so $\varepsilon'\in
K_{5}$, which we excluded. So on every slice of mass $2$,
$E_{1}=\sigma(\delta_{z_{0}}-\delta_{z_{0}\varepsilon'})$. Two slices
$S_{c},S_{c'}$ of mass $2$ are translates of each other: the function
$h(z)=E_{1}(z)-E_{1}(zc^{-1}c')$ on $S_{c}$ is split along $(K_{p})_{p\ne
p_{1}}$ by (S1), has mass at most $4$ and satisfies
$h(z\varepsilon')=-h(z)$, while a nonzero split function of mass at most $4$
is $\pm\mathbf 1_{L}$, which is not of this kind.

If $p_{1}=13$, the twelve slices all have mass $2$, so $E_{1}$ is
$K$-invariant, and $E_{1}=\sigma(\mathbf 1_{L}-\mathbf 1_{L\varepsilon'})$ for
a line $L$ along $K$, with $L\ne L\varepsilon'$: a non-cancelling line sum.

If $p_{1}=11$, then $e^{*}=385$ and $H'=K_{5}\times K_{7}$, of orders $4$
and $6$. The ten slices carry mass at most $24$, at least $2$ each, and
masses at least $3$ on the others; so at least six slices have mass $2$, and
the others come in pairs $\pm c$ with $\mu(c)\le4$. For such a $c$, the
difference of $E_{1}$ on $S_{c}$ with the translate of a slice of mass $2$ is
split, nonzero and of mass at most $\mu(c)+2\le6$, so by \Cref{lem:grid} it
is $\pm\mathbf 1_{L'}$ for a line $L'$ along $K_{5}$ or $K_{7}$. The line
$L'$ contains at most one of the two support points
$z_{0},z_{0}\varepsilon'$ of the translate, as $\varepsilon'$ lies in neither
$K_{5}$ nor $K_{7}$; so $\mu(c)\ge|L'|\ge4$. Hence $\mu(c)=4$, $L'$ is a line along
$K_{5}$ through one of the two points, exactly two slices $S_{\pm c_{1}}$ have
mass $4$, the total is $24$, and $N_{e}=6$. In coordinates $(a,b)\in
(\ZZ/11)^{\times}\times(\ZZ/35)^{\times}$ the support of $E_{1}$ lies in
$\{b=\pm b_{0}\}\cup\{a=\pm c_{1},\ b\in\pm b_{0}K_{5}\}$, where $K_{5}$ is
the group of units congruent to $1$ modulo $7$. For $t=(a,b)$ in this set,
$2t$ lies outside it: $2\not\equiv\pm1$ modulo $35$ and modulo $11$, and
$\pm2,\pm\frac12\notin K_{5}$. By \Cref{lem:signs}(ii) there is then no
top, against $N_{e}=6$.
\end{proof}

\begin{proposition}\label{prop:twoprimes}
Let $e^{*}=pq$ with primes $5\le p<q$. Write $E_{1}(y,x)$ for
$(y,x)\in(\ZZ/p)^{\times}\times(\ZZ/q)^{\times}=U^{*}$. Then
$E_{1}(y,x)=\alpha(x)+\beta(y)$ with odd functions $\alpha,\beta$, either both
with values in $\ZZ$ or both with values in $\frac12+\ZZ$, and
\[
  \|E_{1}\|=\sum_{y,x}\max\bigl(|\alpha(x)|,|\beta(y)|\bigr).
\]
If $e^{*}\notin\{35,55\}$, or if $\alpha$ or $\beta$ is zero, then $E_{1}$ is
a non-cancelling line sum.
\end{proposition}

\begin{proof}
Here $H^{*}=U^{*}$, $K_{p}=(\ZZ/p)^{\times}\times\{1\}$ and
$K_{q}=\{1\}\times(\ZZ/q)^{\times}$, so a split function is a sum of a
function of $x$ and a function of $y$. Fixing $x_{1}$, put
$\alpha_{1}(x)=E_{1}(y,x)-E_{1}(y,x_{1})$, which does not depend on $y$, and
$\beta_{1}(y)=E_{1}(y,x_{1})$; both are integral, and $E_{1}=\alpha_{1}(x)+
\beta_{1}(y)$. Oddness gives $\alpha_{1}(x)+\alpha_{1}(-x)=-\beta_{1}(y)-
\beta_{1}(-y)$, a constant $k\in\ZZ$; so $\alpha=\alpha_{1}-k/2$ and
$\beta=\beta_{1}+k/2$ are odd, with values in $\ZZ$ if $k$ is even and in
$\frac12+\ZZ$ if $k$ is odd. Pairing $(y,x)$ with $(-y,x)$,
$|\alpha(x)+\beta(y)|+|\alpha(x)-\beta(y)|=2\max(|\alpha(x)|,|\beta(y)|)$,
which gives the formula for the mass.

\emph{Half-integral values.} Then $\|E_{1}\|\ge\varphi(e^{*})/2$, so
$\varphi(e^{*})\le48$ and $e^{*}\in\{35,55,65\}$. Let $e^{*}=65$. Equality
holds, so $|\alpha|=|\beta|=\frac12$ everywhere, $N_{e}=6$, $N_{e^{*}}=0$, and
$E_{1}=-1$ exactly on $N=N_{y}\times N_{x}$, where $N_{y}\subseteq
(\ZZ/5)^{\times}$ and $N_{x}\subseteq(\ZZ/13)^{\times}$ are the sets where
$\beta$ and $\alpha$ are negative, of sizes $2$ and $6$. By
\Cref{lem:signs}(ii),(iv), $N$ is the disjoint union of the set $T$ of the six
tops and of $-2T$, and $2T\subseteq-N$. So
$T\subseteq(N_{y}\cap2N_{y})\times(N_{x}\cap(-\tfrac12)N_{x})$, using
$-\frac12\equiv2$ modulo $5$. The set $N_{y}$ contains one element of each
pair $\pm y$, and in $(\ZZ/5)^{\times}=\langle2\rangle$ each of the four such
sets meets its double in one point. So $|T|=6$ forces
$N_{x}=-2N_{x}$; but $-2$ generates $(\ZZ/13)^{\times}$, so $N_{x}$ would be
empty or everything.

\emph{Integral values, $\alpha\ne0\ne\beta$.} As $\alpha$ and $\beta$ are odd,
each is nonzero at two points at least, and $\max(|\alpha|,|\beta|)\ge1$ on at
least $2(p-1)+2(q-1)-4$ points. So $p+q\le16$, and $e^{*}\in\{35,55\}$.

\emph{Integral values, $\alpha=0$ or $\beta=0$.} Then $E_{1}$ is invariant
under $K_{p}$ or under $K_{q}$, and it is a non-cancelling line sum of lines
along that subgroup, each line repeated $|E_{1}|$ times.
\end{proof}

\subsection{The tops \texorpdfstring{$35$}{35} and \texorpdfstring{$55$}{55}}

\begin{lemma}[Carry]\label{lem:carry}
Let $70\mid M$, let $B$ be a multiset in the subgroup of order $70$ of
$\ZZ/M$, and let $A=B+\{y_{1},y_{2}\}$ be balanced, with $y_{1},y_{2}\ne0$
and $s=y_{1}+y_{2}$ of order $35$. If $35$ does not divide $f=\ord(y_{1})$,
then $f\mid70$.
\end{lemma}

\begin{proof}
For a unit $t$,
$\langle ty_{1}\rangle+\langle ty_{2}\rangle=\langle ts\rangle+\delta(t)$
with $\delta(t)=1$ if $\langle ty_{1}\rangle>\langle ts\rangle$ and
$\delta(t)=0$ otherwise. This sum equals $|A|/2-\sum_{b\in B}\langle
tb\rangle$, which depends only on $t$ modulo $70$, and so does $\langle
ts\rangle$; hence so does $\delta(t)$. Suppose $f\nmid70$, and let
$g_{0}=\gcd(f,70)$ and $g_{1}=\gcd(f,35)$, a divisor of $35$ other than $35$.
Write $y_{1}=(M/f)z_{1}$. As $t$ runs over the units congruent to $t'$ modulo
$70$, $\langle ty_{1}\rangle$ runs over the set $P$ of the numbers $z/f$ with
$z$ a unit modulo $f$ and $z\equiv t'z_{1}$ modulo $g_{0}$, by the Chinese
remainder theorem; $P$ depends only on $t'$ modulo $g_{0}$. By the above, $P$
lies on one side of $\langle t's\rangle$, and this holds for every $t'$ in the
class modulo $g_{0}$. Writing $t's=(M/35)w$, the numbers $w/35$ run over the
units $w$ modulo $35$ in a fixed class modulo $g_{1}$, which cut $[0,1)$
into intervals of length at most $14/35$ (if $g_{1}=7$; at most $10/35$ if
$g_{1}=5$, and $3/35$ if $g_{1}=1$). So $P$ lies in an interval of length less
than $1/2$. But $P$ contains two points at distance at least $1/2$:
\begin{itemize}
\item if $f$ has a prime factor $l\notin\{2,5,7\}$, then $l\ge11$; with
$l^{a}\parallel f$ and $z'/f\in P$, $P$ contains the points
$z'/f+i/l^{a}$ for all $i$ outside one class modulo $l$, and these do not lie
in an interval of length less than $1-2/l\ge9/11$;
\item otherwise $4$, $25$ or $49$ divides $f$, $f$ and $g_{0}$ have the same
prime factors, so every $z\equiv t'z_{1}$ modulo $g_{0}$ is a unit modulo $f$,
and $P$ consists of $h=f/g_{0}\ge2$ points spaced by $1/h$, spanning
$(h-1)/h\ge1/2$.
\end{itemize}
This is a contradiction.
\end{proof}

\begin{proposition}\label{prop:3555}
Suppose $e^{*}\in\{35,55\}$ and $E_{1}$ is not a non-cancelling line sum.
Then $M=2e^{*}$.
\end{proposition}

\begin{proof}
Let $n=N_{e}+N_{e^{*}}$, so $\|E_{1}\|=4N_{e}+2N_{e^{*}}\le4n$. By
\Cref{prop:twoprimes}, $\alpha$ and $\beta$ are both nonzero or both
half-integral.

Let $e^{*}=55$. A half-integral $E_{1}$ has mass at least $\varphi(55)/2=20$,
and an integral one with $\alpha,\beta\ne0$ at least $2\cdot4+2\cdot10-4=24$;
so $n\ge5$, and $A$ has at most one low element, which is minus the sum of
the others. All elements of $A$ then lie in the subgroup of order $110$, and
as $A$ generates $\ZZ/M$, $M=110$.

Let $e^{*}=35$, with $y\in(\ZZ/5)^{\times}$ and $x\in(\ZZ/7)^{\times}$. If
$n\ge5$, the same argument gives $M=70$. A half-integral $E_{1}$ has mass at
least $12$ and an integral one with $\alpha,\beta\ne0$ at least $16$; so
$n\ge3$.

\emph{$n=3$.} Then $\|E_{1}\|=12$, $(N_{e},N_{e^{*}})=(3,0)$,
$|\alpha|=|\beta|=\frac12$ everywhere, and $E_{1}=-1$ exactly on
$N=N_{y}\times N_{x}$, with $|N_{y}|=2$ and $|N_{x}|=3$. As in the proof of
\Cref{prop:twoprimes}, $N$ is the disjoint union of $T$ and $-2T$, $|T|=3$, and
$T\subseteq(N_{y}\cap2N_{y})\times(N_{x}\cap3N_{x})$, using
$-\frac12\equiv3$ modulo $7$. Since $|N_{y}\cap2N_{y}|=1$, this forces
$N_{x}=3N_{x}$, impossible as $3$ generates $(\ZZ/7)^{\times}$.

\emph{$n=4$, $\alpha,\beta$ integral and nonzero.} The mass is $16$, so
$(N_{e},N_{e^{*}})=(4,0)$, and equality in the count above gives
$\alpha=s(\delta_{x_{0}}-\delta_{-x_{0}})$,
$\beta=s'(\delta_{y_{0}}-\delta_{-y_{0}})$ with signs $s,s'$. Let $x_{-}$ and
$y_{-}$ be the points where $\alpha$ and $\beta$ equal $-1$, and $x_{+}$,
$y_{+}$ those where they equal $1$. A top $t=(y,x)$ has $E_{1}(t)\le-1$ and
$E_{1}(2t)\ge1$. Since $2\not\equiv\pm1$ modulo $5$ and modulo $7$, at most
one of $x,2x$ lies in $\{\pm x_{0}\}$, and at most one of $y,2y$ in
$\{\pm y_{0}\}$. If $\alpha(x)=-1$, then $\alpha(2x)=0$, so $\beta(2y)=1$,
$\beta(y)=0$ and $t=(y_{+}/2,x_{-})$. If $\alpha(x)=0$, then $\beta(y)=-1$,
$\beta(2y)=0$, $\alpha(2x)=1$ and $t=(y_{-},x_{+}/2)$. The value $\alpha(x)=1$
is impossible. So there are at most two tops, each with $|E_{1}(t)|=1$ and
hence $F(t)=1$ (\Cref{lem:signs}), against $N_{e}=4$.

\emph{$n=4$, half-integral.} A value $|\alpha|\ge\frac32$ or
$|\beta|\ge\frac32$ adds at least $8$ to the mass, which would exceed $16$; so
$|\alpha|=|\beta|=\frac12$, the mass is $12$, $(N_{e},N_{e^{*}})=(2,2)$, and
$A$ has two low elements $y_{1},y_{2}$. The negative set $N=N_{y}\times N_{x}$
is the disjoint union of $-W$, $T$ and $-2T$, where $W$ is the set of the $w$
with $y_{w}\in A$ (\Cref{lem:signs}(iv)). As in the case $n=3$,
$T\subseteq(N_{y}\cap2N_{y})\times(N_{x}\cap3N_{x})$ with
$|N_{y}\cap2N_{y}|=1$, so $N_{y}=\{c,-2c\}$ for some $c$, and $|N_{x}\cap
3N_{x}|\ge2$. In $(\ZZ/7)^{\times}=\langle g\rangle$, $g=3$, an odd set of
three elements meets its product with $g$ in two points exactly when it is
$\{g^{k},g^{k+1},g^{k+2}\}$, and otherwise in none. Hence
\[
  T=\{(c,g^{k+1}),(c,g^{k+2})\},\quad -2T=\{(-2c,g^{k}),(-2c,g^{k+1})\},
  \quad -W=\{(c,g^{k}),(-2c,g^{k+2})\},
\]
using $-2\equiv g^{5}$ modulo $7$. The four elements $x_{t}$, $t\in T$, and
$y_{w}$, $w\in W$, lie in the subgroup of order $70$. In
$\ZZ/70=\ZZ/2\times\ZZ/5\times\ZZ/7$, $x_{t}$ has coordinates $(1,t)$ and
$y_{w}$ has $(0,w)$, so their sum has $2$-coordinate $0$, $5$-coordinate
$c+c-c+2c=3c\ne0$ and $7$-coordinate $g^{k+1}+g^{k+2}-g^{k}-g^{k+2}=2g^{k}
\ne0$. So $s=y_{1}+y_{2}$, which is minus this sum, has order $35$. The
elements $y_{i}$ are low, so $35\nmid\ord(y_{i})$: an order divisible by $35$
has odd part $35$ and is $35$ or $70$. By \Cref{lem:carry}, applied with $B$
the four elements above, $\ord(y_{i})\mid70$. So $A$ lies in the subgroup of
order $70$, and $M=70$.
\end{proof}

\begin{theorem}\label{thm:twoparallel}
Let $M$ be even and prime to $3$, and let $A$ be a sextuple at level $M$ that
contains no pair, is not split and generates $\ZZ/M$. Suppose that $A$ has a
maximal odd part $e^{*}$ whose largest order is $2e^{*}$. Then $A$ has a chain
of moves that succeeds, unless $M\in\{70,110\}$.
\end{theorem}

\begin{proof}
If $A$ has a move of \Cref{lem:elementary} or \Cref{lem:topfibres}, we are
done. Otherwise $e^{*}>1$ is not prime (\Cref{lem:eprime}), and $E_{1}$ is a
non-cancelling line sum by
\Cref{prop:nonsquarefree,prop:threeprimes,prop:twoprimes}, except in the case
of \Cref{prop:3555}, where $M\in\{70,110\}$. In the first case
\Cref{lem:linesum} gives a direct or lowering chain.
\end{proof}
